Pooled ranking gains also arise without grouping (Table~\ref{tab:score-controls}).
In Study A, shrinkage at $\lambda=0.25$ raises pooled AUROC to $0.804$
at $D=300$ and $0.843$ at $D=600$, while preserving the respective raw
within-edge temporal AUROCs $0.707$ and $0.767$. Weights $0.5$ and
$0.75$ also improve pooled ranking. Every positive shrinkage weight preserves
within-edge comparisons, including ties, across the A--C predictions.
Equation~\eqref{eq:score-shrink} attributes these gains to cross-edge reordering.

Window residuals improve the hard-emission comparison. At $D=300$, residual
weights $0.25$, $0.5$ and $0.75$ yield temporal gains over raw of $+0.010$
$[0.001,0.020]$, $+0.009$ $[0.003,0.015]$ and $+0.006$
$[0.003,0.008]$, respectively, compared with hard emission's $-0.069$.
At $D=600$, weight $0.25$ yields $+0.057$ $[0.050,0.064]$, compared
with hard's $+0.038$. These additional analyses report all fixed weights
with pointwise intervals and remain outside the primary confirmatory family.

Changing the rejection output explains much of the gated loss in Study B.
At $D=600$, gains with raw fallback are $+0.013$ $[0.003,0.024]$,
$+0.013$ $[0.004,0.024]$ and $+0.012$ $[0.003,0.022]$ for the
three schedules, compared with static-fallback gains near $-0.21$.
Both rules accept exactly the same candidates. At $D=300$, the corresponding
raw-fallback means are $-0.006$, $0$ and $-0.009$.
Equation~\eqref{eq:fallback-decomposition} separates losses on rejected
repeats from the quality of the accepted candidates.
